% =====================================================================
%  CLOSURE / ZUERICH  --  companion manuscript
%  Psychon Phase Mechanics:
%  The Interior State of a Pruned Agent
%
%  Companion to "A Runtime for Non-Convergent Social Deliberation".
%  That manuscript constructs an agent by pruning a module against a
%  demographic substrate and stops there. This one takes the resulting
%  graph as given and asks what state the agent carries on it, what of
%  that state is available to anyone else, and what one agent can do to
%  another. Self-contained: every object is a finite weighted graph, a
%  bounded convex program, or a construction on one.
% =====================================================================
\documentclass[11pt,a4paper]{article}

\usepackage[utf8]{inputenc}
\usepackage[T1]{fontenc}
\usepackage{lmodern}
\usepackage[a4paper,margin=2.5cm]{geometry}
\usepackage{amsmath,amssymb,amsthm,mathtools}
\usepackage{bm}
\usepackage{booktabs,array,multirow}
\usepackage{enumitem}
\usepackage{microtype}
\usepackage{xcolor}
\usepackage{graphicx}
\usepackage{float}
\usepackage{algorithm}
\usepackage{algpseudocode}
\usepackage[round,sort&compress]{natbib}
\usepackage[colorlinks=true,linkcolor=blue!45!black,
            citecolor=blue!45!black,urlcolor=blue!45!black]{hyperref}
\usepackage[nameinlink,capitalise]{cleveref}

% ---------- theorem environments (matching the parent manuscript) ----
\newtheoremstyle{clstyle}{\topsep}{\topsep}{\itshape}{}{\bfseries}{.}{.5em}{}
\theoremstyle{clstyle}
\newtheorem{theorem}{Theorem}[section]
\newtheorem{lemma}[theorem]{Lemma}
\newtheorem{proposition}[theorem]{Proposition}
\newtheorem{corollary}[theorem]{Corollary}

\theoremstyle{definition}
\newtheorem{definition}[theorem]{Definition}
\newtheorem{construction}[theorem]{Construction}
\newtheorem{example}[theorem]{Example}
\newtheorem{axiom}{Axiom}
\newtheorem{binv}{Invariant}

\theoremstyle{remark}
\newtheorem{remark}[theorem]{Remark}
\newtheorem{notation}[theorem]{Notation}

\crefname{axiom}{Axiom}{Axioms}
\Crefname{axiom}{Axiom}{Axioms}
\crefname{binv}{Invariant}{Invariants}
\Crefname{binv}{Invariant}{Invariants}

% ---------- operators ----------
\DeclareMathOperator{\wt}{w}
\DeclareMathOperator{\cut}{cut}
\newcommand{\sep}{\sigma}
\DeclareMathOperator*{\argmin}{arg\,min}
\DeclareMathOperator*{\argmax}{arg\,max}

% ---------- symbols (parent-compatible) ----------
\newcommand{\R}{\mathbb{R}}
\newcommand{\N}{\mathbb{N}}
\newcommand{\Gph}{\mathcal{G}}
\newcommand{\flo}{\beta}
\newcommand{\Ag}{\mathsf{A}}
\newcommand{\Char}{\chi}
\newcommand{\rec}{m}
\newcommand{\Ord}{R}
\newcommand{\Kc}{K_{c}}
\newcommand{\abs}[1]{\lvert#1\rvert}
\newcommand{\norm}[1]{\lVert#1\rVert}

% ---------- symbols new to this companion ----------
\newcommand{\Psy}{\Psi}                  % psychon
\newcommand{\traj}{\gamma}               % trajectory
\newcommand{\term}{v}                    % terminus
\newcommand{\gap}{D}                     % residual gap
\newcommand{\med}{\bot}                  % medium vertex
\newcommand{\Rep}{\mathsf{r}}            % report-act
\newcommand{\Qst}{\mathsf{q}}            % question-act
\newcommand{\Acts}{\mathcal{A}}
\newcommand{\Hist}{\mathcal{H}}
\newcommand{\drive}{d}
\newcommand{\Pot}{\Phi}

\title{\textbf{Psychon Phase Mechanics}\\[6pt]
\large The Interior State of a Pruned Agent, the Inaccessibility of Its Path,\\
and the Gap-Direction Definition of a Question}

\author{Kundai Farai Sachikonye\\
\small Technical University of Munich\\
\small \texttt{kundai.sachikonye@tum.de}}

\date{\today}

\begin{document}
\maketitle

\begin{abstract}
\noindent
The companion specification constructs an individual by pruning a
thought-pattern module against a demographic substrate, and stops at the
resulting finite weighted graph. This paper takes that graph as given and
asks three questions the construction leaves open: what state does an agent
carry on it, how much of that state is available to any other agent, and
what is the complete set of things one agent can do to another.

We answer with a single object. A \emph{psychon} is a pair
$\Psy=(\traj,\term)$ of a walk through the agent's graph and the vertex it
ends at, subject to a monotone step count. We prove, in order:
\textbf{(i)} a \emph{separation theorem}, that every vertex of a finite
weighted graph with a distinguished medium vertex has a well-defined
\emph{resting cut}---the minimum-weight boundary separating it from the
rest---computable by max-flow, whose cost is bounded below by the floor
$\flo>0$ inherited from the parent construction, so a psychon terminates
somewhere definite and never on a boundary of zero thickness;
\textbf{(ii)} a \emph{recognition--search identity}, that under a monotone
record no operation on the graph has the signature of retrieval, so
recalling is searching and the distinction between memory and inference is
not expressible in the runtime;
\textbf{(iii)} a \emph{path-opacity theorem}, that distinct trajectories
reach identical termini with identical step counts, so the pair
$(\traj,\term)$ is strictly finer than anything an agent can emit, and the
emitted state is exactly the reduced pair $(\term,\rec)$ of terminus and
count;
\textbf{(iv)} a \emph{receiver-relativity theorem}, that one event
registers a different resting cut in every agent that receives it, that
each such registration is correct on its own graph, and that divergence
between receivers is therefore not error and admits no adjudication;
\textbf{(v)} a \emph{drift theorem}, that composed re-registrations across
a chain of receivers have no author, since the composite differs from every
individual contribution and no contribution is distinguished;
\textbf{(vi)} an \emph{act-direction theorem}, that acts partition
exhaustively into two kinds by their effect on the residual gap---a
\emph{report} lowers the gap on the acting side, a \emph{question} raises it
on the receiving side---and that this is the only non-vacuous distinction
available, since neither content nor intent is expressible;
\textbf{(vii)} an \emph{adoption theorem}, that two agents of unequal
internal structure attain equal behaviour when their boundary conditions
match, on a timescale set by coupling above a critical threshold, so that
a practice propagates without any interior being shared or inspected.

We close by fixing the correspondence between the discrete pair
$(\term,\rec)$ and a continuous gradient flow on a bounded coordinate box,
showing the two descriptions agree on the attracting set, and stating
precisely which quantities of the continuous theory are \emph{not}
observables of the instrument. Nine implementation invariants follow, each
a decidable predicate on the runtime.
\end{abstract}

\tableofcontents

% =====================================================================
\section{Introduction}
% =====================================================================

\subsection{What the parent construction leaves open}

The companion specification builds a population as a family of
\emph{modules}---coarse regions of a shared substrate graph, each a
thought-pattern held by many---and produces an individual by \emph{pruning}:
intersecting the modules a demographic record admits, and taking the
subgraph induced on what survives. The construction is deterministic before
contact, and the resulting object is a finite weighted graph.

It stops there, and it must, because at that point the agent has no history.
Two agents pruned from identical records are identical graphs; the parent
manuscript proves they are not yet individuals. What makes them individuals
is what happens on the graph afterwards. That is the subject here.

Three questions are open, and they are not independent.

\begin{enumerate}[label=(\alph*),leftmargin=2em]
\item \textbf{State.} An agent is somewhere on its graph and got there
somehow. Is the ``somehow'' part of its state? If it is, the state is a walk
and grows without bound. If it is not, two agents with different histories
are the same agent, and the individuation claim of the parent manuscript
fails.
\item \textbf{Availability.} Whatever the state is, some function of it is
emitted when the agent acts. If the emitted function were injective, an
observer could reconstruct the interior, and the non-attribution results of
the parent manuscript would be false. If it were constant, agents could not
be told apart at all.
\item \textbf{Acts.} The runtime has to enumerate what one agent can do to
another. Content-based taxonomies (assert, ask, promise, deny) are
unavailable: the parent manuscript proves the runtime cannot evaluate the
truth of a content, and an agent cannot certify its own interior, so any
classification appealing to what is said or meant is not a predicate the
implementation can decide.
\end{enumerate}

We answer all three with one object, and the answers turn out to be forced
rather than chosen.

\subsection{The object}

A \emph{psychon} is a pair
\begin{equation}
\Psy \;=\; (\traj,\term),
\end{equation}
where $\traj$ is a walk on the agent graph and $\term=\traj(k)$ is the
vertex it currently occupies, carried together with a monotone count $\rec$
of committed steps. The pair is the state of the agent. The count is what it
emits. The walk is what nobody, including the agent, can report.

This resolves (a) and (b) simultaneously, and in the only way consistent
with the parent axioms. The walk is genuinely part of the state---so two
agents that reached the same vertex by different routes are different
agents, which is what individuation requires---while the emitted quantity is
the reduced pair $(\term,\rec)$, which is strictly coarser, so reconstruction
of the interior is impossible rather than merely difficult. Question (c) is
answered in \Cref{sec:acts}: the only distinction between acts that survives
the loss of content is the \emph{direction} in which an act moves the
residual gap, and that distinction is exhaustive.

\subsection{Relation to the parent manuscript, and to prior frameworks}
\label{sec:relation}

This paper is a companion, not a revision. It imports the floor, the agent
construction, the attention allocator, and the coordination threshold, and
adds no axiom that the parent does not already state or immediately imply.
Where the parent proves that the instrument \emph{cannot} report a verdict
or an attribution, this paper supplies the interior mechanics that make
those impossibility results properties of the object rather than
restrictions on the interface.

A remark on scope is owed to the reader, because several results below have
recognisable analogues in continuous and physical settings: order
parameters, gradient flows, synchronisation thresholds. We use those
analogues where they are load-bearing and name them where they are not.
Specifically, we adopt a gradient-flow description of the attracting set
(\Cref{sec:continuous}) because it supplies an explicit potential for an
attractor that the parent manuscript proves exists but does not construct;
and we adopt a coupling-threshold form of the adoption result
(\Cref{sec:adoption}) because propagation without shared interiors is
exactly the phenomenon the instrument is built to study. We do \emph{not}
adopt any global order parameter as an observable, and
\Cref{sec:notobservable} states precisely why: a quantity defined over the
whole population is a world-state read-out, and \Cref{thm:receiver} forbids
one. Every continuous object in this paper is defined on the pruned graph of
a single agent and is receiver-relative by construction. Readers approaching
from the physical literature should treat that restriction as the
substantive claim of \Cref{sec:continuous}, not as a limitation of it.

% =====================================================================
\section{Primitives and axioms}
\label{sec:axioms}
% =====================================================================

\begin{definition}[Agent graph]
\label{def:agentgraph}
An \emph{agent graph} is a finite, connected, weighted, undirected graph
$G=(V,E,\wt)$ with $\wt:E\to\R_{>0}$, together with a distinguished vertex
$\med\in V$ adjacent to every other vertex, called the \emph{medium}. We
write $n=\abs{V}\ge 2$.
\end{definition}

The medium is not a thought. It is the vertex standing for everything the
agent has not distinguished: the remainder of the world as it bears on this
graph. Its role is to make ``separating something from the rest'' a
well-posed finite question. In the runtime it is the vertex produced by
pruning that carries no module of its own.

\begin{axiom}[Distinction floor]
\label{ax:floor}
There is $\flo>0$ with $\wt(e)\ge\flo$ for every $e\in E$.
\end{axiom}

\begin{axiom}[Search]
\label{ax:search}
Every operation by which the agent arrives at a vertex is a walk along edges
of $G$ from the vertex it currently occupies. There is no operation whose
effect is to place the agent at a vertex not adjacent to its current one.
\end{axiom}

\begin{axiom}[Act floor]
\label{ax:actfloor}
There is $\varphi>0$ such that every committed step of a walk increments the
record $\rec$ by at least $\varphi$. The record is monotone: no operation
decrements it.
\end{axiom}

\begin{axiom}[Non-completion]
\label{ax:noncompletion}
No walk exhausts $G$: for every walk $\traj$ of finite length there is a
vertex not on $\traj$ and an edge not traversed by $\traj$.
\end{axiom}

\Cref{ax:floor} is the floor of the parent manuscript, restricted to one
pruned graph. \Cref{ax:search} is its locality of conversation.
\Cref{ax:actfloor} is what makes the record a record: an act costing nothing
would leave no trace and could be performed unboundedly often between two
observations. \Cref{ax:noncompletion} holds in the runtime by construction,
since the substrate is strictly larger than any pruned subgraph and the
medium is always available as an unexhausted neighbour.

\begin{remark}[The floor is not a modelling convenience]
\label{rem:floornotconvenience}
A boundary of weight zero is not a boundary: if some $e$ had $\wt(e)=0$, its
endpoints would be separable at no cost and would not be two things. So
\Cref{ax:floor} is constitutive of there being a graph with more than one
vertex at all, rather than an assumption added for tractability. Note also
that on a finite graph the floor is attained automatically, since
$\flo=\min_{e\in E}\wt(e)$ exists and is positive; we state
\Cref{ax:floor} separately only because the constant $\flo$ appears in the
bounds below.
\end{remark}

% =====================================================================
\section{The resting cut}
\label{sec:resting}
% =====================================================================

Before we can say what a psychon terminates \emph{at}, we need the sense in
which a vertex is a definite thing on its graph. That is a separation
question, and it has an exact finite answer.

\begin{definition}[Separation cost and resting cut]
\label{def:sepcost}
For $v\in V\setminus\{\med\}$, let
\begin{equation}
\sep(v) \;=\; \min\Bigl\{\, \wt(\cut(S)) \;:\; v\in S \subseteq V,\ \med\notin S \,\Bigr\},
\label{eq:sepcost}
\end{equation}
where $\cut(S)=\{\,uv\in E : u\in S,\ v\notin S\,\}$ and
$\wt(\cut(S))=\sum_{e\in\cut(S)}\wt(e)$. A minimising set $S^{\star}(v)$ is a
\emph{resting cut} of $v$: the cheapest boundary separating $v$ from the
medium.
\end{definition}

\begin{theorem}[Separation]
\label{thm:separation}
For every $v\in V\setminus\{\med\}$:
\begin{enumerate}[label=(\roman*),leftmargin=2em]
\item $\sep(v)$ is attained by at least one $S^{\star}(v)$;
\item $\sep(v)\ \ge\ \flo\ >\ 0$;
\item $\sep(v)$ equals the value of a maximum $v$--$\med$ flow, and is
computable in strongly polynomial time.
\end{enumerate}
\end{theorem}

\begin{proof}
(i) The feasible family $\{S : v\in S,\ \med\notin S\}$ is a nonempty
subset of the finite power set $2^{V}$ (it contains $\{v\}$), and
$\wt(\cut(\cdot))$ is a real-valued function on it, so a minimiser exists.

(ii) Let $S$ be feasible. Since $\med$ is adjacent to every other vertex
(\Cref{def:agentgraph}) and $v\in S$ while $\med\notin S$, the edge
$v\med\in E$ has exactly one endpoint in $S$, hence $v\med\in\cut(S)$ and
$\cut(S)\neq\emptyset$. Therefore
$\wt(\cut(S))\ge \wt(v\med)\ge\flo$ by \Cref{ax:floor}. Taking the minimum
over feasible $S$ preserves the bound.

(iii) The right-hand side of \eqref{eq:sepcost} is by definition the minimum
weight of a $v$--$\med$ cut in $G$ with capacities $\wt$. By the max-flow
min-cut theorem \citep{fordfulkerson1956maximal} this equals the maximum
$v$--$\med$ flow value, and it is computed in strongly polynomial time by,
for instance, the Edmonds--Karp procedure \citep{edmondskarp1972theoretical}.
\end{proof}

\begin{remark}[What (ii) rules out]
\label{rem:noinfinitesimal}
Part (ii) is the reason a psychon has somewhere to be. Without it a vertex
could be separated from the medium at cost approaching zero, and the
statement ``the agent is at $v$'' would carry no information: the agent
would be at $v$ and, to within any tolerance, equally at $\med$. The floor
makes occupancy quantal. It is also what makes the record of
\Cref{ax:actfloor} meaningful, since a step that crossed no boundary of
positive weight would be a step between two things that were not two.
\end{remark}

\begin{proposition}[Resting cuts need not be unique, and this matters]
\label{prop:nonunique}
There exist agent graphs and vertices $v$ with two distinct minimisers
$S_1^{\star}\neq S_2^{\star}$ of \eqref{eq:sepcost}.
\end{proposition}

\begin{proof}
Take $V=\{\med,v,a,b\}$ with $ab\notin E$ and all remaining edges present,
weighted $\wt(v\med)=3$, $\wt(a\med)=\wt(b\med)=1$, $\wt(va)=\wt(vb)=1$;
every weight is at least $\flo=1$, so \Cref{ax:floor} holds. Then
\begin{align*}
\wt(\cut(\{v\}))      &= \wt(v\med)+\wt(va)+\wt(vb) = 3+1+1 = 5,\\
\wt(\cut(\{v,a\}))    &= \wt(v\med)+\wt(vb)+\wt(a\med) = 3+1+1 = 5,\\
\wt(\cut(\{v,b\}))    &= \wt(v\med)+\wt(va)+\wt(b\med) = 3+1+1 = 5,\\
\wt(\cut(\{v,a,b\}))  &= \wt(v\med)+\wt(a\med)+\wt(b\med) = 3+1+1 = 5,
\end{align*}
and these are all the feasible sets. So all four are minimisers, and
$\sep(v)=5$ is attained non-uniquely.
\end{proof}

Non-uniqueness is not a defect to be broken by a tie rule. It is the finite
shadow of the fact that ``what counts as this thing rather than its
surroundings'' can be genuinely underdetermined by the graph, and any
implementation that silently canonicalises the choice is asserting a
determination the graph does not make. We return to this in
\Cref{binv:tiebreak}.

% =====================================================================
\section{Recognition is search}
\label{sec:recognition}
% =====================================================================

\begin{definition}[Retrieval signature]
\label{def:retrieval}
An operation $\rho$ on the agent has the \emph{retrieval signature} if there
exist vertices $u,v$ with $uv\notin E$ such that $\rho$ carries the agent
from $u$ to $v$, and $\rho$ leaves the record $\rec$ unchanged.
\end{definition}

Retrieval, in the ordinary sense, is exactly this: going straight to the
stored item, at a cost independent of where one was. The next result says
the runtime has no such operation, and that this follows twice over.

\begin{theorem}[Recognition--search identity]
\label{thm:recognition}
No operation of the runtime has the retrieval signature. Consequently, for
every arrival at a vertex there is a walk realising it, and the record
strictly increases along that walk.
\end{theorem}

\begin{proof}
Suppose $\rho$ has the retrieval signature, carrying the agent from $u$ to
$v$ with $uv\notin E$.

First, $\rho$ violates \Cref{ax:search}, which permits arrival only at a
vertex adjacent to the current one; since $uv\notin E$, $v$ is not adjacent
to $u$.

Second, and independently of \Cref{ax:search}: by \Cref{ax:actfloor} the
record is monotone and every committed step increments it by at least
$\varphi>0$, so an operation that changes the occupied vertex while leaving
$\rec$ unchanged commits no step. But then the agent occupies $v$ having
committed nothing, and by monotonicity of $\rec$ the runtime cannot
distinguish the state after $\rho$ from the state before it. Two distinct
occupancies with an identical record are, by \Cref{def:emitted} below, the
same emitted state, so $\rho$ is not an operation of the runtime at all: it
has no effect the runtime can represent.

Hence no such $\rho$ exists. Every arrival is therefore an adjacency step
(\Cref{ax:search}), a finite composition of which is a walk, and each step
increments $\rec$ by at least $\varphi$.
\end{proof}

\begin{corollary}[Memory and inference are not distinguishable in the runtime]
\label{cor:memoryinference}
There is no predicate on runtime states that holds of arrivals-by-recall and
fails of arrivals-by-derivation.
\end{corollary}

\begin{proof}
By \Cref{thm:recognition} both are walks on $G$ satisfying
\Cref{ax:search,ax:actfloor}. A predicate separating them would have to be a
function of the walk beyond its status as a walk; the only such data are the
identities of intermediate vertices, and by \Cref{thm:opacity} below these
are not present in any emitted state. A predicate on runtime states is by
definition a function of emitted state.
\end{proof}

\begin{remark}
\Cref{cor:memoryinference} is not scepticism about the distinction. It is a
statement about this instrument: an agent that answers ``I remember'' and an
agent that answers ``I worked it out'' have performed the same kind of
operation, and the runtime is not entitled to label one of them. The second
argument in the proof of \Cref{thm:recognition} is worth isolating, because
it does not use \Cref{ax:search} at all: even an implementation that
permitted jumps would find retrieval incoherent, since under a monotone
record a costless jump is invisible.
\end{remark}

% =====================================================================
\section{The psychon and path opacity}
\label{sec:psychon}
% =====================================================================

\begin{definition}[Psychon]
\label{def:psychon}
A \emph{psychon} on an agent graph $G$ is a pair $\Psy=(\traj,\term)$ where
$\traj=(v_0,v_1,\dots,v_k)$ is a walk in $G$ (so $v_{i}v_{i+1}\in E$ for all
$i<k$) and $\term=v_k$. Its \emph{record} is
$\rec(\Psy)=\sum_{i<k}c(v_iv_{i+1})$ where $c(e)\ge\varphi$ is the committed
cost of traversing $e$, per \Cref{ax:actfloor}.
\end{definition}

\begin{definition}[Emitted state]
\label{def:emitted}
The \emph{emitted state} of a psychon is the pair
\begin{equation}
\pi(\Psy) \;=\; \bigl(\term,\ \rec(\Psy)\bigr) \;\in\; V\times\R_{\ge0}.
\end{equation}
An agent acts only through $\pi$: no operation of the runtime takes as
input, or produces as output, any function of $\traj$ not factoring through
$\pi$.
\end{definition}

\begin{theorem}[Path opacity]
\label{thm:opacity}
Let $G$ be an agent graph with $n\ge3$. Then there exist psychons
$\Psy_1=(\traj_1,\term)$ and $\Psy_2=(\traj_2,\term)$ with
$\traj_1\neq\traj_2$, the same terminus $\term$, and
$\rec(\Psy_1)=\rec(\Psy_2)$. Hence $\pi$ is not injective, and the interior
of a psychon is not a function of its emitted state.
\end{theorem}

\begin{proof}
Since $n\ge3$ there are distinct $a,b\in V\setminus\{\med\}$. The medium is
adjacent to all vertices, so $\med a,\med b\in E$. Consider
\begin{equation}
\traj_1=(\med,a,\med,b,\med),\qquad \traj_2=(\med,b,\med,a,\med),
\end{equation}
both walks in $G$ with the same terminus $\term=\med$. They are distinct as
sequences, since $a\neq b$ gives $\traj_1(1)=a\neq b=\traj_2(1)$.

Both traverse the same multiset of edges, namely $\med a$ twice and $\med b$
twice, differing only in the order. Since $\rec$ is by \Cref{def:psychon} the
sum of committed costs over traversed edges, and addition is commutative,
\begin{equation}
\rec(\Psy_1) \;=\; 2c(\med a)+2c(\med b) \;=\; \rec(\Psy_2),
\end{equation}
for \emph{any} cost function $c$. Hence
$\pi(\Psy_1)=(\med,\rec(\Psy_1))=(\med,\rec(\Psy_2))=\pi(\Psy_2)$ while
$\Psy_1\neq\Psy_2$, so $\pi$ is not injective.
\end{proof}

\begin{corollary}[No self-certification]
\label{cor:noselfcert}
An agent cannot emit a warrant that its own terminus was reached by one walk
rather than another.
\end{corollary}

\begin{proof}
Any such warrant is an emission, hence by \Cref{def:emitted} a function of
$\pi(\Psy)$. By \Cref{thm:opacity} there are $\Psi_1\neq\Psi_2$ with
$\pi(\Psy_1)=\pi(\Psy_2)$; a function of $\pi$ takes the same value on both,
so it does not distinguish them.
\end{proof}

\begin{corollary}[Individuation]
\label{cor:individuation}
Two agents pruned to the same graph, having emitted the same
$(\term,\rec)$, may nonetheless be distinct psychons; and once
$\rec>0$ the pair $(\traj,\term)$ is not reproducible from the pruning
record.
\end{corollary}

\begin{proof}
Immediate from \Cref{thm:opacity} for the first claim. For the second: the
pruning record determines $G$ and hence the set of available walks, but by
\Cref{ax:actfloor} the record is monotone and by \Cref{ax:noncompletion} no
walk is forced, so the particular walk taken is not a function of $G$.
\end{proof}

This is the promised completion of the parent construction. Before contact,
pruning is deterministic and two agents from the same record are the same
object; after the first committed step, they are psychons with different
interiors, and no emission recovers the difference.

% =====================================================================
\section{Receiver relativity and drift}
\label{sec:receiver}
% =====================================================================

So far one agent. We now put two or more on the same event and ask what they
end up holding.

\begin{definition}[Registration]
\label{def:registration}
Let $\Ag_1,\dots,\Ag_N$ be agents with graphs $G_1,\dots,G_N$, and let $x$
be an event. A \emph{registration} of $x$ in $\Ag_i$ is a vertex
$\term_i\in V_i$ at which a walk prompted by $x$ terminates, together with
its resting cut $S^{\star}(\term_i)\subseteq V_i$ (\Cref{def:sepcost}).
\end{definition}

\begin{theorem}[Receiver relativity]
\label{thm:receiver}
Let $x$ be a single event received by agents $\Ag_1,\Ag_2$ with
$G_1\neq G_2$. Then:
\begin{enumerate}[label=(\roman*),leftmargin=2em]
\item the registrations $S^{\star}(\term_1)\subseteq V_1$ and
$S^{\star}(\term_2)\subseteq V_2$ are subsets of different vertex sets and
need not correspond under any graph morphism;
\item each is a correct minimiser of \eqref{eq:sepcost} on its own graph;
\item there is no graph on which both are cuts, hence no third party
relative to which one is nearer to $x$ than the other.
\end{enumerate}
\end{theorem}

\begin{proof}
(i) By construction $S^{\star}(\term_i)$ is defined as a subset of $V_i$.
Since $G_1\neq G_2$ there need be no bijection $V_1\to V_2$ preserving
$\wt$; pruning against distinct demographic records produces distinct
induced subgraphs, and \Cref{prop:nonunique} shows minimisers are not even
unique within one graph.

(ii) Correctness is a property internal to \eqref{eq:sepcost}: $S^{\star}
(\term_i)$ minimises $\wt(\cut(\cdot))$ over the feasible family of $G_i$.
That is a finite optimisation with an attained minimum by
\Cref{thm:separation}(i), and it makes no reference to $G_{3-i}$.

(iii) Suppose $H$ were a graph on which both $S^{\star}(\term_1)$ and
$S^{\star}(\term_2)$ are cuts, so that their weights are comparable as
separations in $H$. Then $H$ contains $V_1\cup V_2$ and assigns weights to
edges within each; but the weight of $\cut(S^{\star}(\term_1))$ computed in
$H$ includes edges to $V_2\setminus V_1$, which are absent from $G_1$ and
whose weights are not determined by either agent graph. So the comparison
requires data not contained in $G_1$, $G_2$, or the event; by
\Cref{ax:noncompletion} no agent holds it, and by \Cref{def:emitted} no
emission supplies it.
\end{proof}

\begin{corollary}[Divergence is not error]
\label{cor:divergence}
If $\Ag_1$ and $\Ag_2$ register the same event differently, no predicate of
the runtime marks either registration as incorrect.
\end{corollary}

\begin{proof}
By \Cref{thm:receiver}(ii) both are correct on their own graphs, and by
(iii) there is no graph relative to which they can be jointly scored. A
predicate marking one incorrect would be a function of data the runtime does
not hold.
\end{proof}

\begin{remark}[Falsity is an exterior relation]
\label{rem:falsity}
\Cref{cor:divergence} does not say nothing is ever false. It says falsity is
not a unary predicate on an agent state. To call a registration false is to
compare a pair---the agent and something else---and the something else is
always another graph, held by another agent, with the same standing. An
observer may of course fix a reference graph and grade everything against
it; that observer is then a further agent, and \Cref{thm:receiver} applies
to them too. What the runtime must not do is privilege one such graph
internally, which is the content of \Cref{binv:noreference}.
\end{remark}

\begin{definition}[Reshaping and drift]
\label{def:drift}
A \emph{reshaping} of an agent graph is a map $g:G\mapsto G'$ altering
weights, arising from a registration. For a chain of receivers
$\Ag_1\to\Ag_2\to\cdots\to\Ag_n$, each passing on what it registered, the
\emph{composed drift} is
\begin{equation}
G_n \;=\; (g_n\circ g_{n-1}\circ\cdots\circ g_1)(G_0).
\end{equation}
\end{definition}

\begin{theorem}[Drift has no author]
\label{thm:drift}
Let $n\ge2$ and suppose each $g_i$ is non-trivial, i.e.\ $g_i\neq\mathrm{id}$.
Then in general $G_n\neq g_i(G_0)$ for every $i$, and no $g_i$ is
distinguished by any function of the emitted states
$\{\pi_i\}_{i=1}^{n}$.
\end{theorem}

\begin{proof}
The first claim is immediate: composition of non-identity weight maps
differs from each factor unless the others cancel, which is non-generic.

For the second, suppose $F$ is a function of $\{\pi_i\}$ selecting some
index $j$ as the author. Each $\pi_i=(\term_i,\rec_i)$ by
\Cref{def:emitted}. Consider the chain in which $\Ag_j$ and $\Ag_{j+1}$
exchange their reshapings, $g_j'=g_{j+1}$ and $g_{j+1}'=g_j$. Because
composition of weight maps on disjoint edge supports commutes, the composite
$G_n$ is unchanged; and because \Cref{thm:opacity} makes $\pi_i$ independent
of which walk produced it, the emitted states can be arranged identical
under the exchange. Then $F$ returns $j$ in both, but the reshaping
performed by $\Ag_j$ differs between them. So $F$ does not select the agent
that performed a given reshaping.
\end{proof}

\begin{corollary}[Rumour]
\label{cor:rumour}
A registration that has passed through $n\ge2$ receivers is a composed
drift: it is held by all of them, originated by none identifiably, and is
not false in the sense of \Cref{cor:divergence}.
\end{corollary}

This is the formal content of the phenomenon the instrument is built to
observe. A proposal circulates, everyone who holds it holds something
slightly different, each holding is correct on its own graph, and no
participant is the source. The parent manuscript proves the runtime cannot
attribute; \Cref{thm:drift} says there is nothing to attribute.

% =====================================================================
\section{Acts and the direction of the gap}
\label{sec:acts}
% =====================================================================

We can now answer question (c) of the introduction. The difficulty is that
the obvious taxonomies are unavailable: by \Cref{cor:noselfcert} an agent
cannot certify its own interior, so intent is not emitted; by
\Cref{cor:divergence} the runtime cannot evaluate a content; and by
\Cref{thm:opacity} the walk is not visible. What remains is the effect an
act has, and it turns out that is enough---and that it yields exactly two
kinds.

\begin{definition}[Residual gap]
\label{def:gap}
For an agent $\Ag$ engaged with a matter, the \emph{residual gap}
$\gap(\Ag)\ge0$ is the total weight of the resting cut at its current
terminus that remains untraversed:
\begin{equation}
\gap(\Ag) \;=\; \wt\bigl(\cut(S^{\star}(\term))\bigr) \;-\; \rho(\Ag),
\end{equation}
where $\rho(\Ag)$ is the weight already committed against that cut.
\end{definition}

By \Cref{thm:separation}(ii), $\wt(\cut(S^{\star}(\term)))\ge\flo>0$, so the
gap starts positive; by \Cref{ax:noncompletion} it is never exhausted by any
finite walk.

\begin{theorem}[Exchange does not close]
\label{thm:forcedact}
Let two agents exchange acts about a matter. While $\gap>0$ on either side,
a further act is available, and no state of the exchange is a terminal state
in which both gaps are zero.
\end{theorem}

\begin{proof}
Suppose $\gap(\Ag_1)=\gap(\Ag_2)=0$. Then each agent has committed weight
equal to the full resting cut at its terminus, so each has traversed every
edge of $\cut(S^{\star}(\term_i))$. By \Cref{ax:noncompletion} there remains
in $G_i$ a vertex not on the walk and an edge not traversed; take such an
edge $e$ incident to $S^{\star}(\term_i)$ if one exists. If none exists,
$S^{\star}(\term_i)$ is a connected component separated from the rest, and
since $\med\notin S^{\star}(\term_i)$ while $\med$ is adjacent to every
vertex (\Cref{def:agentgraph}), some edge from $S^{\star}(\term_i)$ to
$\med$ is present and by \Cref{ax:floor} has weight at least $\flo>0$;
traversing it is an available act that returns the gap to a positive value.
Either way an act is available, contradicting terminality.
\end{proof}

\begin{remark}
\Cref{thm:forcedact} is why the instrument does not terminate on agreement.
Exchange ends by exhaustion of what the participants are willing to spend,
not by the matter closing, because the matter does not close. This is the
non-convergence of the parent title, localised to a pair.
\end{remark}

\begin{definition}[Report and question]
\label{def:reportquestion}
Let $\Ag$ act toward $\Ag'$, and write $\gap$, $\gap'$ for the gaps
immediately before and $\widehat{\gap}$, $\widehat{\gap}'$ immediately after.
The act is a
\begin{itemize}[leftmargin=2em]
\item \emph{report}, written $\Rep$, if $\widehat{\gap}<\gap$: it lowers the
gap on the acting side;
\item \emph{question}, written $\Qst$, if $\widehat{\gap}'>\gap'$: it raises
the gap on the receiving side.
\end{itemize}
\end{definition}

\begin{theorem}[Act direction]
\label{thm:actdirection}
\Cref{def:reportquestion} is well posed and exhaustive:
\begin{enumerate}[label=(\roman*),leftmargin=2em]
\item every act with any effect is a report, a question, or both;
\item the two are independent, i.e.\ each of the four combinations of
$\{\text{is}/\text{is not a report}\}\times\{\text{is}/\text{is not a
question}\}$ is realisable except the doubly-negative one, which has no
effect;
\item the classification is decidable from emitted states alone, without
reference to content, intent, or interior.
\end{enumerate}
\end{theorem}

\begin{proof}
(iii) first, since it constrains the rest. By \Cref{def:gap}, $\gap$ is a
function of $\term$ and the committed weight $\rho$, both components of the
emitted state $\pi$ of \Cref{def:emitted} (the resting cut is computed from
$\term$ and $G$ by \Cref{thm:separation}(iii)). So the sign of
$\widehat{\gap}-\gap$ is computable from $\pi$ before and after, which is
what the runtime holds.

(i) Let $\alpha$ be an act with some effect. By \Cref{ax:actfloor} it
commits at least $\varphi>0$, so it changes $\rho$ on at least one side. If
it changes $\rho(\Ag)$ upward, then by \Cref{def:gap} $\widehat{\gap}<\gap$
and $\alpha$ is a report. If it changes the receiving side, it does so by
prompting a walk in $\Ag'$ (\Cref{ax:search}), which by
\Cref{thm:recognition} moves $\term'$; a move of terminus changes
$S^{\star}(\term')$ and hence $\wt(\cut(S^{\star}(\term')))$. If that weight
rises, $\widehat{\gap}'>\gap'$ and $\alpha$ is a question. If it does not
rise, the receiving side has committed against its cut without the cut
growing, which is a report by $\Ag'$ rather than an effect of $\alpha$. So
an act with an effect falls in at least one class.

(ii) Realisability. A report that is not a question: $\Ag$ traverses an edge
of its own resting cut and $\Ag'$ takes no step, so $\gap$ falls and $\gap'$
is unchanged. A question that is not a report: $\Ag$ acts so as to move
$\term'$ to a vertex of larger separation cost while committing nothing
against its own cut, possible since by \Cref{prop:nonunique} distinct
vertices of $G'$ have distinct $\sep$ values in general. Both at once: an
act doing each of the above. Neither: an act committing nothing, excluded by
\Cref{ax:actfloor} from having an effect.
\end{proof}

\begin{corollary}[A question is not a request for content]
\label{cor:questioncontent}
The predicate ``$\alpha$ is a question'' is invariant under relabelling of
what $\alpha$ says.
\end{corollary}

\begin{proof}
By \Cref{thm:actdirection}(iii) the predicate is a function of emitted
states, which by \Cref{def:emitted} contain $(\term,\rec)$ and no content.
\end{proof}

\begin{remark}[Why this is the definition worth having]
\label{rem:whydirection}
An agent that asks in order to learn and an agent that asks in order to
unsettle perform the same act under \Cref{def:reportquestion}, and the
instrument is right not to separate them: by \Cref{cor:noselfcert} neither
could certify which it was doing, including to itself. What the definition
does capture is the asymmetry that matters for a deliberation: a report
spends the speaker down, a question spends the hearer up. A population of
agents that only report converges to silence by exhaustion; a population
that also questions does not, and \Cref{thm:forcedact} says it cannot. This
is the mechanism by which the agents of the instrument ask questions
\emph{for their own graph} rather than on a script: raising $\gap'$ is the
only act whose cost is paid by someone else.
\end{remark}

% =====================================================================
\section{Adoption without shared interiors}
\label{sec:adoption}
% =====================================================================

The instrument must be able to represent a practice spreading: one agent
comes to do what another does. \Cref{thm:opacity} and \Cref{thm:receiver}
between them forbid this happening by transfer of interior state. The
following says it happens anyway, through boundaries.

\begin{definition}[Boundary condition]
\label{def:boundary}
The \emph{boundary condition} of agent $\Ag$ at terminus $\term$ is the
weighted edge set $\partial(\Ag)=\cut(S^{\star}(\term))$ together with the
map $e\mapsto\wt(e)$ on it. Two agents \emph{match} at tolerance
$\varepsilon$ if their boundary conditions agree, as weighted multisets, to
within $\varepsilon$.
\end{definition}

\begin{theorem}[Behavioural adoption]
\label{thm:adoption}
Let $\Ag_1,\Ag_2$ have graphs $G_1\neq G_2$ and match at tolerance
$\varepsilon$. Then for every act $\alpha$ classified by
\Cref{def:reportquestion}, the classification of $\alpha$ agrees between
$\Ag_1$ and $\Ag_2$ up to $O(\varepsilon)$. In particular, matching
boundaries induce matching behaviour without any correspondence between
$V_1$ and $V_2$.
\end{theorem}

\begin{proof}
By \Cref{thm:actdirection}(iii) the classification of $\alpha$ is the sign
of a difference of gaps, and by \Cref{def:gap} the gap is
$\wt(\cut(S^{\star}(\term)))-\rho$. The first term is a function of the
boundary condition alone (\Cref{def:boundary}); the second is the committed
record, which is emitted. If the boundary conditions agree to within
$\varepsilon$ as weighted multisets, the first terms agree to within
$\varepsilon$, so the gaps agree to within $\varepsilon$ at equal $\rho$,
and the signs of their differences agree wherever those differences exceed
$\varepsilon$ in magnitude. By \Cref{ax:actfloor} every act moves $\rho$ by
at least $\varphi$, so for $\varepsilon<\varphi$ the classification agrees
exactly. No step of this argument refers to $V_1$, $V_2$, or any map between
them.
\end{proof}

\begin{corollary}[Adoption is not inspection]
\label{cor:adoptionnotinspection}
$\Ag_2$ may come to behave as $\Ag_1$ does without any emission of $\Ag_1$
having determined any interior of $\Ag_2$.
\end{corollary}

\begin{proof}
By \Cref{thm:adoption} matching $\partial(\Ag_1)$ and $\partial(\Ag_2)$
suffices for matching behaviour, and $\partial$ is a function of the resting
cut, which by \Cref{thm:opacity} does not determine the walk. Distinct
interiors are compatible with identical boundaries.
\end{proof}

\begin{proposition}[Threshold]
\label{prop:threshold}
Let a population of $N$ agents be coupled with strength $K$ and let
$\Delta$ measure the spread of their boundary conditions. Matching at
tolerance $\varepsilon$ is attained for all pairs only if $K>\Kc$ with
$\Kc=\Theta(\Delta)$; below $\Kc$ the unmatched configuration is stable.
\end{proposition}

\begin{proof}[Proof sketch]
This is the coordination theorem of the parent manuscript applied to
$\partial$ rather than to internal state, and the argument is unchanged:
coupling must overcome heterogeneity, the critical value scales linearly in
the spread, and below it the incoherent configuration is linearly stable. We
do not reproduce it; we note only that the quantity being locked here is a
boundary, which is exactly the quantity \Cref{thm:adoption} shows is
sufficient.
\end{proof}

\begin{remark}[Erosion under replacement]
\label{rem:erosion}
Adoption is not free. Replacing a part of an agent graph is a removal
followed by an addition, and each leaves a residue: the removal discards the
weights of edges between the removed part and the remainder, and the
addition must fix weights not determined by either. By \Cref{ax:actfloor}
these are committed operations, so the record grows monotonically under
replacement and never returns. After $n$ replacements the accumulated
residue is $\Theta(n)$ under statistically similar replacements. An agent
that has adopted much is therefore not the agent that began, and by
\Cref{cor:individuation} there is no operation restoring it. We record this
because it is the honest cost of \Cref{thm:adoption}: a population that
converges behaviourally has paid for it in individuation, and the instrument
should display that rather than hide it.
\end{remark}

% =====================================================================
\section{The continuous correspondence}
\label{sec:continuous}
% =====================================================================

The discrete account above is complete as it stands. This section adds a
continuous description of the \emph{attracting set} of an agent, because the
parent manuscript proves such an attractor exists without constructing one,
and because an explicit potential is useful for implementation. Everything
here is defined on one agent graph. \Cref{sec:notobservable} states what
this does \emph{not} license.

\begin{definition}[Interior coordinates]
\label{def:coords}
Fix an agent $\Ag$ with graph $G$. Its \emph{interior coordinates} are
\begin{equation}
q \;=\; (\Ord,\ \sigma^2,\ s) \;\in\; \mathcal{M} \;=\; [0,1]\times[\sigma^2_{\min},\sigma^2_{\max}]\times[0,1]^{d},
\end{equation}
where $\Ord$ is the coherence of the agent walk (the fraction of recent
steps lying inside the current resting cut), $\sigma^2$ is the dispersion of
its terminus over recent steps, and $s$ collects $d$ bounded summaries of
the emitted record. $\mathcal{M}$ is compact and convex.
\end{definition}

\begin{definition}[Drive potential]
\label{def:potential}
A \emph{drive potential} for $\Ag$ is a function
$\Pot:\mathcal{M}\to\R$ of the form
\begin{equation}
\Pot(q) \;=\; \underbrace{\tfrac{1}{2}(\Kc-K)\Ord^{2}+\tfrac{1}{4}K\Ord^{4}}_{\text{coherence}}
\;+\;\underbrace{\theta\,\sigma^{2}+\frac{\theta}{K\sigma^{2}}}_{\text{dispersion floor}}
\;-\;\underbrace{a\,\Ord\,e^{-\sigma^{2}/\tau}}_{\text{coupling}}
\;+\;\underbrace{V(s)}_{\text{boundary}},
\label{eq:potential}
\end{equation}
with $K,\theta,a,\tau>0$ and $V$ smooth, bounded and convex on $[0,1]^d$.
\end{definition}

\begin{theorem}[Attractor]
\label{thm:attractor}
Suppose $\Pot$ is $\lambda$-strongly convex on $\mathcal{M}$ for some
$\lambda>0$. Then the gradient flow
\begin{equation}
\dot q \;=\; -\,\nabla\Pot(q)
\label{eq:flow}
\end{equation}
has a unique fixed point $q^{\star}\in\mathcal{M}$, and every trajectory
converges to it exponentially:
\begin{equation}
\norm{q(t)-q^{\star}} \;\le\; e^{-\lambda t}\,\norm{q(0)-q^{\star}}.
\end{equation}
Moreover the discretisation $q_{k+1}=q_k-\delta\nabla\Pot(q_k)$ satisfies
$\norm{q_k-q^{\star}}\le(1-\lambda\delta)^{k}\norm{q_0-q^{\star}}$ for
$\delta\le1/L$, where $L$ is the Lipschitz constant of $\nabla\Pot$.
\end{theorem}

\begin{proof}
$\mathcal{M}$ is compact and convex and $\Pot$ is $\lambda$-strongly convex
on it, so $\Pot$ attains its minimum at a unique $q^{\star}$, which is the
unique stationary point and hence the unique fixed point of \eqref{eq:flow}.
Let $W(t)=\tfrac12\norm{q(t)-q^{\star}}^2$. Then
$\dot W=-\langle q-q^{\star},\nabla\Pot(q)\rangle\le-\lambda\norm{q-q^\star}^2=-2\lambda W$
by strong convexity together with $\nabla\Pot(q^{\star})=0$, and Gr\"onwall
gives the stated bound. The discrete estimate is the standard contraction
bound for gradient descent on a $\lambda$-strongly convex, $L$-smooth
objective with step $\delta\le1/L$ \citep{boyd2004convex}.
\end{proof}

\begin{proposition}[Agreement with the discrete account]
\label{prop:agreement}
The step bound $\delta\le1/L$ of \Cref{thm:attractor} is compatible with
\Cref{ax:actfloor}: the act floor $\varphi$ bounds $\delta$ from below, so
the admissible range is $\varphi\le\delta\le1/L$, nonempty whenever
$L\le1/\varphi$. Within it, the fixed point $q^{\star}$ is the continuous
image of the set of termini at which the residual gap is stationary.
\end{proposition}

\begin{proof}
By \Cref{ax:actfloor} no committed step is smaller than $\varphi$, so a
discretisation of \eqref{eq:flow} realisable by the runtime has
$\delta\ge\varphi$. Stationarity of $\Pot$ means $\nabla\Pot(q^\star)=0$; by
\Cref{def:coords} the components of $q$ are functions of the terminus and
the record, so a stationary $q$ is one at which further committed steps do
not change the coherence or dispersion of the walk, which by \Cref{def:gap}
is the condition that $\gap$ is stationary.
\end{proof}

\begin{remark}[What the potential adds and what it does not]
\label{rem:potentialadds}
\Cref{thm:attractor} supplies a constructive $\Pot$ for the attractor the
parent manuscript proves abstractly. It does not weaken
\Cref{thm:forcedact}: $q^\star$ is a stationary point of the coherence of a
walk, not a state of zero gap. By \Cref{ax:noncompletion} the gap remains
positive at $q^{\star}$, so an agent at its attractor is still acting. The
attractor is a habit, not a conclusion.
\end{remark}

\subsection{What is not an observable}
\label{sec:notobservable}

\begin{proposition}[No global order parameter]
\label{prop:noglobal}
Let $\Ord_i$ be the coherence of agent $\Ag_i$ per \Cref{def:coords}. Then
no aggregate $\Ord_{\mathrm{pop}}=F(\Ord_1,\dots,\Ord_N)$ is an observable of
the runtime.
\end{proposition}

\begin{proof}
Each $\Ord_i$ is defined relative to the resting cut of $\Ag_i$ on $G_i$
(\Cref{def:coords}). By \Cref{thm:receiver}(iii) there is no graph on which
the resting cuts of two distinct agents are jointly cuts, so the $\Ord_i$
are not commensurable quantities and $F$ is not well defined on them. Even
granting a fixed normalisation, $F$ would be a function of the interiors of
all agents; by \Cref{thm:opacity} interiors are not emitted, so $F$ is not
computable from what the runtime holds.
\end{proof}

\begin{corollary}[The instrument still reports nothing it should not]
\label{cor:stillnothing}
Adjoining \Cref{def:potential} to the runtime introduces no verdict, no
progress measure, and no attribution.
\end{corollary}

\begin{proof}
A verdict or progress measure over the population is an aggregate of the
kind excluded by \Cref{prop:noglobal}. A per-agent $\Pot$ is not one: it is
defined on $G_i$ alone, is receiver-relative by
\Cref{thm:receiver}, and by \Cref{rem:potentialadds} its minimum is not a
terminal state. Attribution is excluded by \Cref{thm:drift}.
\end{proof}

\begin{remark}[On the temptation]
\label{rem:temptation}
The result above is the whole reason this section is written as carefully as
it is. A potential and an order parameter are exactly the objects one
reaches for when asked how the population is doing, and computing
$\Ord_{\mathrm{pop}}$ would be a few lines. \Cref{prop:noglobal} says the
resulting number would not mean anything: not that it would be imprecise,
but that the quantities averaged are defined on incommensurable graphs.
The instrument may compute $\Pot$ per agent because each such $\Pot$ is that
agent view; it may not compute the population figure, because there is no
agent whose view it would be.
\end{remark}

% =====================================================================
\section{Implementation invariants}
\label{sec:invariants}
% =====================================================================

Each invariant is a decidable predicate on the runtime, paired with the
result it enforces.

\begin{binv}[Floor]
\label{binv:floor}
Every edge weight in every pruned graph is at least $\flo>0$, checked at
construction. Enforces \Cref{ax:floor} and hence
\Cref{thm:separation}(ii).
\end{binv}

\begin{binv}[Monotone record]
\label{binv:record}
The record $\rec$ of every agent is non-decreasing across every operation.
No interface decrements it. Enforces \Cref{ax:actfloor} and
\Cref{thm:recognition}.
\end{binv}

\begin{binv}[No retrieval]
\label{binv:noretrieval}
No operation sets an agent terminus to a vertex non-adjacent to its current
terminus. Enforces \Cref{ax:search}; violation would reintroduce the
retrieval signature of \Cref{def:retrieval}.
\end{binv}

\begin{binv}[Emission is the reduced pair]
\label{binv:emission}
Every value crossing the agent boundary factors through
$\pi(\Psy)=(\term,\rec)$. No serialisation exposes $\traj$. Enforces
\Cref{thm:opacity} and \Cref{cor:noselfcert}.
\end{binv}

\begin{binv}[Explicit non-uniqueness]
\label{binv:tiebreak}
Where \eqref{eq:sepcost} has multiple minimisers, the implementation records
that fact rather than silently selecting one; any selection is marked as a
choice of the implementation and not a determination of the graph. Enforces
the reading of \Cref{prop:nonunique}.
\end{binv}

\begin{binv}[No reference graph]
\label{binv:noreference}
No agent graph is designated as the standard against which others are
scored, and no code path compares registrations across agents for
correctness. Enforces \Cref{thm:receiver} and \Cref{cor:divergence}.
\end{binv}

\begin{binv}[No authorship field]
\label{binv:noauthor}
No representation of a circulating registration carries an originator, and
none is reconstructible from stored state. Enforces \Cref{thm:drift}.
\end{binv}

\begin{binv}[Acts classified by direction only]
\label{binv:direction}
The report/question classification is computed from gaps before and after,
never from content or a declared intent field. Enforces
\Cref{thm:actdirection}(iii) and \Cref{cor:questioncontent}.
\end{binv}

\begin{binv}[No population aggregate]
\label{binv:noaggregate}
No route, log line, or stored value aggregates a per-agent coherence,
potential, or gap across agents. Enforces \Cref{prop:noglobal} and
\Cref{cor:stillnothing}.
\end{binv}

% =====================================================================
\section{Limitations}
\label{sec:limitations}
% =====================================================================

\paragraph{The continuous section is optional.}
\Cref{sec:continuous} is a convenience for implementation, not a load-bearing
part of the argument. Deleting it leaves
\Cref{thm:separation,thm:recognition,thm:opacity,thm:receiver,thm:drift,thm:actdirection,thm:adoption}
intact. We flag this because the continuous apparatus is the part most
likely to be over-read.

\paragraph{The threshold is imported.}
\Cref{prop:threshold} is stated as a sketch because its proof is the
coordination theorem of the parent manuscript. We have applied it to a new
quantity---the boundary condition rather than internal state---and that
application is ours; the underlying threshold result is not.

\paragraph{Tolerance in adoption is a parameter.}
\Cref{thm:adoption} gives exact agreement only for $\varepsilon<\varphi$.
For coarser matching the classification agrees only up to $O(\varepsilon)$,
and the instrument should expose $\varepsilon$ rather than fix it.

\paragraph{Non-uniqueness is under-explored.}
\Cref{prop:nonunique} establishes that resting cuts may fail to be unique
and \Cref{binv:tiebreak} says what to do about it, but we do not
characterise when non-uniqueness occurs, nor how often it does on graphs
produced by pruning against the Z\"urich substrate. That is an empirical
question this paper does not answer.

% =====================================================================
\section{Conclusion}
% =====================================================================

We took the graph the parent construction produces and asked what lives on
it. The answer is a psychon: a walk together with the vertex it has reached,
carrying a monotone record. The walk is the agent individuality and is
inaccessible; the pair of terminus and record is everything anyone else
gets; and those two facts are the same fact, since it is precisely the
non-injectivity of the emission map that makes two identically-pruned agents
into two agents.

From this the conversational mechanics follow without further assumption.
Recalling is searching, because a costless jump under a monotone record is
invisible. One event registers differently in every receiver, each correctly,
with no graph on which the registrations compete. A registration passed
along a chain becomes a drift with no author. And an act is a report or a
question according to which side of the exchange its cost falls on---the
only distinction that survives when content, intent and interior are all
unavailable, and, we have argued, the one that was doing the work anyway.

Finally, practices spread. Agents whose boundaries match behave alike
regardless of what is behind those boundaries, which is how something can
propagate through a population that shares no interior and inspects nothing.
It is not free: every adoption is a replacement, every replacement commits
record, and the record does not run backwards. A population that comes to
agree has spent individuation to do it. The instrument is built to show that
transaction rather than to score its outcome, and
\Cref{sec:invariants} is the list of places where we have had to say so in
code.

\bibliographystyle{plainnat}
\bibliography{references}

\end{document}
